\chapter[Sermon 26. Of Him That Sits in Throne, and Holds the Book in His Right Hand Sealed with vii. Seals: What That Sealed Book Is.]{Of Him That Sits in Throne, and Holds the Book in His Right Hand Sealed with vii. Seals: What That Sealed Book Is.}
\label{sermon:026}
\markright{Sermon 26. Of Him That Sits in Throne, and Holds the Book in His ...}

And I saw in the right hand of him who sat on the throne a book written inside and on the back, sealed with seven seals. And I saw a mighty angel proclaiming with a loud voice, “Who is worthy to open the book and to loose its seals?” And no one in heaven, nor on earth, nor under the earth, was able to open the book and to look upon it. And I wept greatly because no one was found worthy to open and read the book, nor even to look upon it.

He now proceeds to describe more fully him that sits on the Throne: of whom he had touched certain things before. This section contains considerable weight for our matter. For now he will show that which is principal in this treatise: that all things which are done in the world through God’s providence are most justly and holily governed by Christ. All the saints of God, and all creatures, acknowledging this, for an example to us, do praise and celebrate him that lives forever.

And it shall behoove us to weigh every word, since in each one are great mysteries, and nothing is spoken in vain. And truly, God Almighty sits in a throne. And by sitting is signified not only the power of judging, ruling, and governing, but also a quiet mind (not troubled with any evil affections, after the manner of judges of this world) and great equity in all things. Secondly, a book is seen in the right hand of him that sits, of which book we must speak more at large.

\index{Apocalypse (Book of)}\index{Daniel (Book of)}Here appears an allusion, as there is in many other places of the Scripture, to the princes of this world, who have books of laws, of privileges, of institutes—regarding what has been done and what is to be done—finally, of secrets, of acts, of the condemned, and of citizens, of life and of death. For so is both the book and books assigned to God. Moses says in Exodus 32:33, “Put me out of the book of life,” and so forth. In the Psalms there is much mention of these books of God: in Psalms 56, 69, and 139. In the seventh of Daniel, books are opened, of which mention is also made in the twentieth of the Apocalypse. We read in Malachi 3 of a book of remembrance before God. Therefore, this book of God contains all the counsels of God, all his works and judgments. For we shall hear by and by that all things that are done in the world come out of this book, as it were out of a fountain or wellspring.

And three things are chiefly spoken of this book. First, that it lies not in the Throne, or in the bosom of him that sits, or under the Throne, or that it hangs before or behind the Throne: but it is in the right hand of God. Hereby is signified the operation or power of God, and the same most just and most mighty. For the book is not seen in the left hand. God therefore works, and contains or ministers all his works and judgments most holily. Secondly, that book is written within and without, or on the backside. For in God’s providence and judgments, all things are contained, both good and evil, fortunate and unfortunate, sharp and gentle, sweet and sour, visible and invisible, secret and apart, and all things in general.

Finally, the book is sealed with seven seals. For it is most strongly closed and fastened. For the judgments and works of God are firm, true, just, and such as cannot be withstood. The use of seals among men is diverse, yet it may be considered in two points. First, seals are set because of fidelity, truth, and righteousness. And great deliberation is had in setting the seals. For they are not put to unjust, vain, or false matters. Therefore, seals are tokens of certainty and testimonies of a right. It seems an unworthy thing to speak against sealed writings. By the seals therefore that are set to the book of God is signified that the judgments and works of God are most firm, true, and just, whatever is done by his providence, and are ordained by Christ. It shall therefore be a shame to find fault with the judgments of God, or to speak evil of his works. Again, by seals, secrets are kept, that they be not seen by every man, but by them only to whom they are appointed. The judgments therefore and works of God are for the most part hid, and not open to all men, saving to such as the Lord hath appointed, namely to the faithful and obedient. But there are only seven seals, for in them the fullness of times, and of things to be done in these times throughout the world and church, and of the judgments and mysteries of God are comprehended.

Now therefore, the opening of the book and its unsealing, are nothing else but the revealing of God’s judgments and the declaring or uttering of his most secret counsels. Finally, the most holy and just operation, dispensation, and execution of his will. Nothing in that opening is done against the truth, faith, love, and justice of God.

And with many words, and also most diligently and well is treated here of the opening of the seals, who truly might be thought worthy to open to the church the secret judgments of God, and to execute and minister his holy works: that is to say, to whom the kingdom is given and government of the divine providence. For an Angel, and that not of the common sort, but a strong and worthy one, with a loud voice cries, to make us all attentive, and that we should note diligently, who he is that should both open the book and unloose or undo the seals. And he holds the hearer, beholder, or reader in suspense before he will show him, to the intent truly to commend him to us exceedingly. No man, says he, in the whole universal world, neither among the angels and saints in heaven, nor among earthly men and under the earth, was found who could either open or unseal the book.

\index{John, Saint}\index{Papacy}Let us observe that no one can open the book and its seals except Christ alone. Why then is the administration of things attributed or communicated to saints? No one can reveal to us the counsels and judgments of God, nor can any man govern those judgments and works of God that he works in the world, save only Christ the Lord. Why then are such great benefits sought from saints, and credited to them, if either the sick are restored to health, or if a mortal man receives any other gift or benefit? Many will say, “I received this as a gift of God, but through the meditation, power, and merit of this or that saint, to whom God granted authority to rule over such a disease and might heal those who call upon the name of the saint, or the name of God by the saint.” These things are now refuted by the words of the Lord and Saint John, who say that no one in heaven or on earth is found who could open the book. Yet nevertheless, about the throne sat the twenty-four elders, representing the type of all saints in glory, and not one of them was found who could open the book. Therefore, they are greatly mistaken who attribute the governance of things in the church to the Pope, a most corrupt and filthy man. Only Christ received all power in heaven and on earth, as we shall presently understand more fully.

John wept, for he understood that a weighty matter consisted in the opening of this godly book; and yet he saw no one at all who could either open or unseal it. Neither did he as yet fully understand the matter. And he bore the figure of those who understand not the judgments of God, nor know that all things are through God’s providence holily governed by Christ. For in them nothing else remains but mourning and heaviness. Certainly, without Christ and his opening, whereby he reveals to us the divine mysteries and judgments, no man can rightly judge of the same. For unless we understand the seals to be opened by Christ, and that all things are done by his order, which loved us and gave himself for us, what thing shall be left in us but sighing?

He presented three things, to open, read, and look upon. No living person can open [it], for no person is fit for so great a responsibility, except the Son of God. No person reads or fully understands God’s judgments, but the Son, and to whomever he has revealed as much as any person has. No person can look upon [it], that is to say, can behold the works and judgments of God, but will be offended, unless he is endowed with the Spirit and purified by it. Therefore, we must ask grace from him, that we may understand so much of God’s judgments as will suffice, and may judge them well.

\index{Arethas of Caesarea}Arethas, Bishop of Caesarea, an expositor of this book, says that neither those who lack a body, nor those who are embodied, nor those who have departed, leaving their bodies behind, have received a perfect understanding of divine matters. And immediately afterward: neither was there anyone who could open it, nor even one who could attentively consider it—that is to say, could not carefully observe the judgments of God. And the context of the entire passage sufficiently proves that Saint John speaks here of the judgments truly, but especially of the governance of things. The Lord Jesus be glorified forever. Amen.
